
\subsubsection{Overview}

Our central hypothesis is that certain environment-stage API defects violate documented invariants testable before training begins. We stratify these invariants into three classes requiring distinct validation mechanisms:

\textbf{Structural invariants} (tensor shapes, dtypes, device placement) are detectable at import time through static introspection.

\textbf{Metamorphic invariants} (mathematical properties like probability sums, identity relations) require lightweight runtime evaluation on synthetic inputs.

\textbf{Composition invariants} (cross-library consistency in device placement, dtype propagation) arise from multi-library interactions and require validation at framework boundaries.

\subsubsection{Tier 1: Structural Contracts}

\textbf{Invariants Validated}: Structural contracts encode shape, dtype, and device constraints---the interface properties that function signatures specify but Python's type system cannot enforce.

\textbf{Implementation}: Structural contracts are Python decorators that wrap function definitions:

\begin{verbatim}
@validate_structural(
    expected_input_channels=3,
    expected_dtype=torch.float32
)
def forward(x: torch.Tensor) -> torch.Tensor:
    ...
\end{verbatim}

At import time, decorators intercept the first function call to validate actual arguments against specifications.

\textbf{Design Rationale}: Import-time validation enables fail-fast behavior---researchers discover mismatches immediately upon importing modules rather than during training. Our proof-of-concept demonstrates validation overhead below 0.03 seconds, orders of magnitude faster than training-stage detection.

\subsubsection{Tier 2: Metamorphic Contracts}

\textbf{Invariants Validated}: Metamorphic contracts enforce mathematical properties---softmax outputs must sum to 1.0, dropout must preserve expectation under eval mode, batch normalization must not change distributional statistics during inference.

\textbf{Implementation}: Metamorphic contracts assert input-output relations via lightweight probes:

\begin{verbatim}
@validate_metamorphic(
    property='softmax_sum',
    probe=lambda f, x: torch.allclose(
        f(x).sum(dim=-1), 
        torch.ones(...), 
        atol=1e-5
    )
)
def softmax(x: torch.Tensor, dim: int = -1) -> torch.Tensor:
    ...
\end{verbatim}

Contracts execute probes on synthetic inputs before the first production call.

\textbf{Design Rationale}: By executing probes once at environment-setup on synthetic inputs, we amortize validation cost across the entire training run. If softmax violates probability-sum invariants on synthetic inputs, it will likely violate them on real data.

\textbf{Floating-Point Tolerance}: ML computations involve approximate arithmetic where exact equality fails even for mathematically equivalent expressions. Contracts use \texttt{torch.allclose(atol=1e-5)} for numeric comparisons, with tolerance thresholds derived from IEEE 754 single-precision limits.

\subsubsection{Tier 3: Composition Contracts}

\textbf{Invariants Validated}: Composition contracts validate cross-library consistency---PyTorch tensors passed to framework conversion tools must preserve dtype, shape, and numerical tolerance.

\textbf{Implementation}: Composition contracts intercept cross-library boundaries:

\begin{verbatim}
with validate_composition(
    contracts=['dtype_preserved', 'shape_consistent', 'numerical_tolerance'],
    rtol=1e-5, atol=1e-7
):
    target_model = convert_framework(source_model, inputs)
\end{verbatim}

The validator wraps conversion operations and validates dtype preservation, shape consistency, and numerical equivalence.

\textbf{Design Rationale}: Cross-framework conversion introduces opportunities for dtype corruption, shape mismatches, and numerical drift. Validating these properties at conversion boundaries enables pre-execution defect detection.

\subsubsection{Execution Model}

\begin{verbatim}
Import Time (0-50ms)
├─ Structural contracts: Validate tensor shapes, dtypes, devices
└─ Register metamorphic/composition contracts

Environment Setup (50-500ms)
├─ Metamorphic contracts: Execute probes on synthetic inputs
└─ Composition contracts: Validate library bindings

Training Begins (>500ms)
└─ Contracts dormant; no per-batch overhead
\end{verbatim}

Contracts incur zero per-batch overhead during training---validation occurs once at environment-setup, then contracts become dormant.
